\section{Methodology}
\label{sec:method}

We describe our approach for extracting geometric features from attention circuits and testing whether they predict prompt-level performance.

\subsection{Task and Data}
\label{sec:data}

We use \ssttwo (Stanford Sentiment Treebank, binary) from the GLUE benchmark for sentiment classification (positive/negative).
We sample 200 examples (100 positive, 100 negative) from the validation set with balanced random sampling (seed=42).
We choose \ssttwo because it is the most commonly used dataset for prompt sensitivity studies \citep{sclar2023quantifying, mizrahi2024state} and provides a clean binary classification signal.
Texts are truncated to 200 characters to ensure manageable sequence lengths within \gptsmall's 512-token context window.

\subsection{Prompt Format Variants}
\label{sec:prompts}

We create 20 systematically varied prompt formats, inspired by the format perturbation methodology of \citet{sclar2023quantifying}.
Formats vary along several axes: instruction phrasing (\eg ``Classify the sentiment'' vs.\ ``Is the following review positive or negative?''), separator tokens (newlines, pipes, tabs, colons), structural markup (XML tags, JSON-like, parenthetical), label words (positive/negative, good/bad, great/terrible), and instruction casing (standard, all-caps).
\tabref{tab:formats} in \secref{sec:results} provides the complete list with accuracy values.
This systematic variation isolates formatting effects while preserving semantic content.

\subsection{Performance Measurement}
\label{sec:performance}

We run \gptsmall (12 layers, 12 heads, 768-dimensional hidden states) on all 200 examples across all 20 formats.
For each example, we compare the logit scores for the positive and negative label tokens and assign the predicted label as the one with the higher logit.
Accuracy per format is the fraction of correct predictions across all 200 examples.
We define a binary failure label using the median accuracy (0.603) as the threshold: formats below the median are labeled as failures, and those above as stable.
This yields a balanced 10/10 split.

\subsection{Geometric Feature Extraction}
\label{sec:features}

We use TransformerLens \citep{nanda2022transformerlens} to extract attention patterns from all 12 layers $\times$ 12 heads of \gptsmall.
For computational tractability, we subsample 50 examples from our 200-example set for cache extraction.
For each prompt format and example, we obtain the full attention matrix $\mA^{(l,h)} \in \sR^{T \times T}$ for each layer $l \in \{1,\ldots,12\}$ and head $h \in \{1,\ldots,12\}$, where $T$ is the sequence length.

We compute 23 geometric features organized into six categories:

\para{Entropy features (7).}
We compute the Shannon entropy of each attention head's distribution: $H(\mA^{(l,h)}) = -\sum_{i,j} a_{ij} \log a_{ij}$.
We aggregate across heads and examples to obtain the mean, standard deviation, and maximum entropy per format.
We also compute layer-group entropy (early: layers 1--4, mid: layers 5--8, late: layers 9--12) and the standard deviation of entropy across layers.

\para{Singular value features (4).}
We compute the SVD of each attention matrix and extract: (i) the effective rank $r_{\text{eff}} = \exp(H(\hat{\sigma}))$, where $\hat{\sigma}$ is the normalized singular value distribution; (ii) the condition number $\kappa = \sigma_1/\sigma_{\min}$; (iii) the ratio of the top-3 singular values to the total; and (iv) the entropy of the singular value spectrum.
These features capture the rank structure and concentration of attention patterns \citep{bussman2025beyond}.

\para{Variance features (5).}
We compute the variance of attention weights within each head, then aggregate to obtain the mean, standard deviation, maximum, and mid-layer variance across the format.
The standard deviation of variance across layers captures cross-layer consistency.
High variance indicates complex attention landscapes that may be less robust \citep{park2024geometry}.

\para{Concentration features (2).}
We measure how peaked each attention distribution is using the maximum attention weight per position, then compute the mean and standard deviation across heads.

\para{Inter-head similarity (2).}
We compute pairwise cosine similarity between attention patterns of different heads within each layer, yielding the mean and standard deviation.
This captures head redundancy versus diversity.

\para{Frobenius norm features (2).}
We compute the Frobenius norm $\|\mA^{(l,h)}\|_F$ of each attention matrix, then aggregate to obtain the mean and standard deviation across heads and layers.
The Frobenius norm measures overall attention magnitude, related to the query-weight norm that \citet{bali2026quantifying} found to correlate with head stability.

All features are computed per-head, per-layer, then averaged across the 50 subsampled examples to obtain a single feature vector $\vf \in \sR^{23}$ per prompt format.

\subsection{Statistical Analysis}
\label{sec:analysis}

We compute Spearman rank correlations between each geometric feature and prompt accuracy.
We use a permutation test (5{,}000 permutations) to assess whether the top correlation is significant beyond what would be expected from testing 23 features.
We also compute bootstrap 95\% confidence intervals (10{,}000 resamples) for each correlation.

\subsection{Failure Classification}
\label{sec:classifier}

We train logistic regression classifiers to predict the binary failure label from geometric features.
Given the small sample size ($n = 20$), we use leave-one-out (LOO) cross-validation.
We evaluate classifiers using the top-$k$ features (for $k \in \{2, 3, 5\}$) selected by absolute Spearman correlation on the full dataset.

\para{Baselines.}
We compare against two baselines: (i) \emph{random prediction}, which assigns labels uniformly at random (expected 50\% accuracy); and (ii) \emph{logit confidence}, which uses the mean maximum softmax probability across examples as a proxy for prompt quality \citep{hendrycks2017baseline}.
